\chapter{結論}
\label{chap:conclusion}
本研究では，カメラ画像に基づく経路生成，速度計画，経路追従制御からなる自律走行システムを対象に，ラップタイムを決定している要因を，要素別・上限速度の段階別に定量化した．研究目的の3点に対する結果は次のとおりである．

第一に，ラップタイムを，車両の物理的な限界だけで決まる基準時間と，上限速度の設定・速度計画・経路生成・経路追従・停止に起因する時間差に分ける方法を，実車走行と閉ループシミュレーションに共通の物差しとして構成し，実車のR1（赤コース，上限速度4.0\,m/s，Pure Pursuit法）とB1（青コース，上限速度4.5\,m/s，MPC）に適用した（4章）．R1は，1周73.0\,sのうち基準時間が59.6\%，上限速度の設定値による時間差が29.6\%，信号停止が10.1\%を占め，速度計画と経路生成に由来する時間差は3周の合計で$-0.2$\,sであった．B1は，基準時間が72.1\%，上限速度の設定値による時間差が23.6\%で，速度計画と経路生成に由来する時間差は解析範囲（2.97周）の合計で$+6.65$\,s（ラップタイムの3.7\%，1周あたり約2.2\,s）であり，その約半分（$+3.29$\,s）は分岐での速度抑制による．上限速度の設定値4\,m/sの段階では，短縮の余地の大部分は上限速度の設定値にあり，速度計画と経路生成の設定に由来する時間差は，R1でほぼなく，B1で1周あたり約2.2\,sである．

第二に，走行試験で同定した車両応答と，周回走行の記録から求めた経路生成の誤差を組み込んだ閉ループシミュレーションを構成し，各要素を置き換えた，または外した場合のラップタイムと完走を比較した（5章）．シミュレーションは，R1相当の3周のラップタイムの平均が72.7\,s（実車73.0\,s），B1相当の1，2周目が60.5，59.6\,s（実車60.54，59.75\,s）と，1\%以内の差で実車に近い値になり，内訳の各項の差も1\,s以内であった（追従の項はシミュレーションが0.7〜0.8\,s大きい）．一方，信号停止の行き過ぎの符号（シミュレーションは3回とも停止線の手前，実車は2回が越えた）は再現していない．比較の結果は次のとおりである．完走を左右するのは局所経路計画の合流距離であり，15\,mに固定すると青コースの完走は0/5，赤コースの完走は2/5〜3/5になる．赤コースで合流の実効距離を変えると，最大横偏差は実効距離10\,m以下の条件で0.40〜0.55\,m，約12\,mで1.3〜1.4\,m，15\,mで1.5〜1.6\,mであったが，機構は特定できていない．信号停止を含まない周回では，MPCをPure Pursuit法に置き換えてもラップタイムの差は0.5\,s以内で完走は変わらず，最大横偏差が4〜7倍になる．信号停止を含むR1相当では，MPCは距離誤差ありのとき停止1回あたりの時間差が約15\,s（Pure Pursuit法は約7\,s）となってラップが80〜82\,sになり，3周の完走は4/5であった．分岐での速度抑制は，青コースで1周あたり3.7\,s（5.5\,m/s），4.3\,s（6.0\,m/s）を要する．平滑化の次数とカーブでの認識距離の短縮は，ラップタイムにも完走にも差を生じず，経路生成の雑音をなくしても完走は変わらず，ラップタイムの差は0.4\,s以内であった．信号停止は，実測の距離誤差を与えた検出モデルで，4.0〜7.0\,m/sの5/5で速度0.05\,m/s未満まで停止し，行き過ぎは最大$+0.86$\,mであり，停止専用の減速度の効果は小さかった．

第三に，上限速度の段階ごとの制約要因を，条件式とシミュレーションで示した（5章）．4.0\,m/sでは上限速度の設定値が，5.5\,m/sでは上限速度の設定値に加えて青コースの分岐での速度抑制が時間差を生じ，青コースでは6.0\,m/s以上で上限速度の設定値による時間差が残らず，抑制が主な時間差になる（6.0\,m/sで抑制ありの基準49.3\,s，抑制なし45.1\,s）．条件式は，停止については，実距離で停止専用の減速度が上限まで強められる場合の見積もり（9.7\,m/s）がシミュレーションと矛盾しない．先読みについては，認識距離5〜9\,mの仮定で5.8〜7.5\,m/sを与え，シミュレーションで抑制なしの交差点カーブの横加速度が速度計画の上限を超え始める4.5\,m/sより高く，条件式が予測する認識距離への依存は，シミュレーションでは再現されなかった．なお，6.0\,m/sを超える結果は，試験で到達を確認した最高速度を超える参考値である．

これらの結果の適用範囲は，実車の走行がR1，B1の2本であること，走行可能な最高速度は試験で到達を確認した6.0\,m/sであり，横加速度の上限が仮定値であること，シミュレーションの自己位置が真値で経路生成のモデルが揺らぎだけを含むこと，障害物回避を評価していないことに限られる．また，MPCと信号停止の組合せは実車で確認しておらず，R1相当の構成で3周の完走が4/5となった原因は調べていない．

今後の課題として，第一に，上限速度の設定を6.0\,m/sより上げた試験で，走行可能な最高速度と横加速度の上限を同定すること，第二に，直進での右回りの偏り（従動輪の零点，左右輪の半径差，路面の傾斜）を切り分けること，第三に，経路生成の誤差と検出のモデルは周回走行R1，B1から同定しており，同じ走行で妥当性を確認しているため，別の周回走行で検証すること，第四に，5\,m/s以上の実車走行で，分岐での速度抑制なしの交差点カーブの横加速度と，MPCでの信号停止を確認すること，第五に，合流距離が完走を左右する機構を特定することが挙げられる．
